\section{符号说明}

\begin{center}
\begin{tabular}{>{\centering\arraybackslash}p{0.16\textwidth}p{0.55\textwidth}>{\centering\arraybackslash}p{0.16\textwidth}}
\toprule
符号 & 含义 & 单位 \\
\midrule
\(O\) & 目标圆域中心，即坐标原点 & m \\
\(S_i\) & 第 \(i\) 个检测点坐标 & m \\
\(G_j\) & 第 \(j\) 个干扰源坐标 & m \\
\(\theta_i\) & 检测点 \(S_i\) 测得的示向度 & \(\degree\) \\
\(\varepsilon\) & 示向度最大绝对误差，题设为 \(1\degree\) & \(\degree\) \\
\(\Omega\) & 多次示向扇形的公共定位区域 & --- \\
\(D(\Omega)\) & 定位区域的直径 & m \\
\(R_j\) & 干扰源 \(j\) 的有效接收半径 & m \\
\(T_{\mathrm{all}}\) & 一次测试的定位清除总时间 & s \\
\(N, N_c\) & 干扰源总数与已清除个数 & 个 \\
\bottomrule
\end{tabular}
\end{center}

仅在单个问题中使用的符号于首次出现处定义。
